\section{Related Work}
\label{sec:ch2-relatedwork}

In recent years, the use of machine learning for modelling partial differential equations has become a central topic in scientific machine learning. Data-driven surrogate models have repeatedly been shown to achieve accuracy comparable to classical numerical solvers while offering orders-of-magnitude speedups at inference time. Convolutional, graph-based, and transformer-based architectures have all found success in this setting, with transformer-based \textit{foundation models} attracting particular interest over the last two years due to their ability to generalise across PDE families, parameter regimes, and boundary conditions. In parallel, the recognition that high-dimensional, chaotic systems such as turbulent fluid flow exhibit branching behaviour has motivated a shift from purely deterministic prediction towards probabilistic modelling. This section reviews classical surrogate and operator-learning models, transformer-based PDE foundation models, foundation models developed outside the classical PDE literature but built on related architectural ideas, the datasets that have enabled this progress, work on bifurcative and chaotic dynamical systems, and flow matching as a generative framework for probabilistic modelling.

\subsection{Surrogate models and neural operators} 

Learning-based surrogates for PDEs have existed for close to a decade. Convolutional Neural Networks (CNN) such as the U-Net~\cite{ronneberger2015unet}, originally developed for biomedical image segmentation, remain a popular backbone for physics surrogates because their encoder-decoder structure with skip connections balances local and global feature learning at a comparatively low computational cost. Rather than a fixed-resolution image-to-image mapping in CNNs, a new approach formulates the learning problem as \textit{operator learning}, essentially learning a mapping between infinite-dimensional function spaces. The Fourier Neural Operator (FNO) ~\cite{li2021fno} learns PDE solution operators by parameterising convolution kernels in the frequency domain, which allows it to capture long-range spatial dependencies efficiently and to generalise across different discretisations. DeepONet~\cite{lu2021deeponet} instead uses two separate networks, a branch network encoding the input function and a trunk network encoding the query coordinates, and has since been extended to multiple input functions in MIONet~\cite{jin2022mionet}. Graph neural network-based operators~\cite{sanchezgonzalez2020learning, pfaff2021meshgraphnets} extend surrogate modelling to irregular meshes and particle-based representations, whereas standard convolutional modelling uses static grid structures as input. A further widely used technique is to embed physical constraints directly into the training loss, referred to as Physics Informed Neural Networks (PINNs)~\cite{raissi2019pinn}, which regularise the learned solution to satisfy the conservation laws of the governing PDE, though such models are known to suffer from optimisation pathologies on stiff or highly non-linear systems~\cite{krishnapriyan2021failuremodes}. While convolutional models and neural operators are comparatively cheap to train and are robust on smaller, single-PDE datasets, they generalise poorly to unseen equation forms, parameters, or boundary conditions, motivating the shift towards larger, better generalising models based on transformer architectures discussed next.

\subsection{Transformer models for Fluid Dynamics} 

Transformer architectures, popularised in natural language processing by models such as GPT and LLaMA, combine tokenisation or patching with self-attention layers to capture both local and long-range dependencies. In computer vision, the Vision Transformer (ViT)~\cite{dosovitskiy2020vit} and the hierarchical Swin Transformer~\cite{liu2021swin}, which introduces shifted local windowed attention and progressive downsampling, established transformers as a competitive and highly scalable alternative to convolutional backbones. These architectural trends have since been carried over to PDE modelling, where transformer-based models trained and evaluated across a broad range of PDE types and parameter settings are commonly referred to as \textit{PDE foundation models}. Specifically for fluid dynamics, Multiple Physics Pretraining (MPP)~\cite{mccabe2024mpp} pretrains an autoregressive vision transformer with a novel axial attention mechanism across heterogeneous physics datasets, encouraging the emergence of shared representations across PDE systems. DPOT~\cite{hao2024dpot} instead combines an auto-regressive denoising pretraining objective, in which Gaussian noise is injected into the input frames to mitigate train-test mismatch during rollout, with a Fourier-attention transformer backbone. Where standard transformers use a static patch size, MATEY~\cite{zhang2024matey} proposes multiscale, adaptive patching together with several attention variants to handle spatiotemporal physical systems at varying resolutions. PROSE-FD~\cite{liu2024prosefd}, building on the earlier PROSE-PDE framework~\cite{liu2024prosepde}, additionally encodes the symbolic form of the governing equation alongside the numerical state, fusing a symbolic and a data encoder so that the underlying physics can inform the operator prediction. Poseidon~\cite{herde2024poseidon} introduces the scalable Operator Transformer (scOT), a Swin transformer-based multiscale U-Net, pretrained on compressible Euler and incompressible Navier-Stokes data using a novel all to all prediction method. VICON~\cite{cao2024vicon} takes a different approach, casting multi-physics operator learning as \textit{in-context} learning. This vision-transformer-based in-context operator network infers the underlying operator directly from example trajectories at inference time rather than relying purely on pretrained weights, and is reported to substantially reduce long-term rollout error relative to MPP. Related autoregressive transformer approaches include BCAT~\cite{liu2025bcat}, which uses block-causal attention for high-dimensional fluid states and report efficiency-accuraccy tradeoff of different spatiotemporal attention mechanisms.

Despite this progress, most of these models share a set of limitations that are rarely discussed explicitly. First, pretraining datasets remain comparatively small and are typically drawn from a limited number of simulations. Models are instead validated primarily on their ability to transfer across PDE types, rather than generalising across the far larger space of conditions of a single PDE form. Second, predictions are conditioned on only a short window of preceding states, typically a single timestep, and the training objective is almost always single-step or fixed-horizon prediction. Due to compounding errors in auto-regressive use, these models are inherently unstable when rolled out over long horizons. Due to the combination of limited intra-equation variability during pretraining and the neglect of long-horizon temporal stability, we adress the gap in research by training on a large, heterogenous traing dataset and focus on long-horizon rollout stability in the training and validation

\subsection{Foundation models beyond Fluid Dynamics} 

The same architectural ideas have proven effective beyond the classical PDE-benchmark setting. In plasma physics, Swin-Transformer/U-Net hybrids have been used as neural surrogates for five-dimensional gyrokinetic turbulence simulations~\cite{gyrokinetic2024}. In atmospheric science, ClimaX~\cite{nguyen2023climax} introduces a ViT-based foundation model for weather and climate that uses variable-specific tokenisation and variable aggregation to cope with heterogeneous sets of physical variables across data sources, while Aurora~\cite{bodnar2024aurora}, developed by Microsoft, scales this idea to a large 3D Swin Transformer with Perceiver-based encoders and decoders, pretrained across multiple heterogeneous atmospheric datasets and fine-tuned with LoRA for long-lead-time rollout. Outside the regular-grid setting, the Erwin transformer~\cite{erwin2025} proposes a graph-based hierarchical attention mechanism for physical systems that are naturally represented as point clouds rather than structured grids. 

Taken together, this body of work confirms that hierarchical, windowed-attention transformer backbones transfer well across a wide range of physical domains, which is used as inspiration for our own model.

\subsection{Datasets} 

The development of these models has been enabled by large simulated PDE datasets. For unobstructed flow governed by compressible and incompressible Navier-Stokes and related fluid-dynamics equations, PDEBench~\cite{takamoto2022pdebench}, PDEArena~\cite{gupta2022pdearena} are widely used. PDEGym~\cite{herde2024poseidon} was specifically simulated for Poseidon and focusses on a wide variety of incompressible intitialisation conditions. CFDBench~\cite{luo2023cfdbench} and FlowBench~\cite{tali2024flowbench} focus on obstructed and boundary-driven flows, while MPFBench and BubbleML~\cite{hassan2023bubbleml} target multiphase, bubble, and droplet dynamics. BLASTNet~\cite{chung2023blastnet} and The Well~\cite{ohana2024thewell} provide large-scale collections of multi-physics simulation data spanning several PDE families and physical regimes. The diversity of these sources, in equation form, boundary and initial conditions, grid type, and numerical solver, is what makes cross-dataset generalisation a non-trivial challenge. 

\subsection{Bifurcative/sensitive systems} 

High-dimensional or chaotic systems, of which turbulent Navier-Stokes flow is an example, frequently exhibit branching behaviour in their evolution. Small perturbations to the initial condition can lead to qualitatively different, yet equally physically valid, future trajectories, commonly referred to as bifurcations. Modelling such systems with a deterministic, single-trajectory predictor is flawed as such models can only have a one-to-one mapping. This motivates a probabilistic approach to prediction, in which the model represents a distribution over plausible trajectories rather than committing to a single point estimate, and is the primary reason we compare a deterministic autoregressive model against a probabilistic, flow-matching alternative.

\subsection{Flow matching and generative modelling} 

Compared to the standard deterministic modelling methods above, Flow Matching (FM)~\cite{lipman2022flowmatching} allows 


case~\cite{lipman2024flowmatchingguide}. Rectified Flow~\cite{liu2022rectifiedflow} re-trains the model on its own straightened trajectories to reduce the number of integration steps needed at inference. STFlow~\cite{tenbrinke2025stflow} instead modifies the starting point itself: rather than Gaussian noise, it initialises from a random walk fit to the observed initial frames, so the model only has to learn a correction on top of an already-informed prior, which is reported to reduce transport cost and improve accuracy specifically on trajectories prone to bifurcation. This combination of cheap, simulation-free training and a flexible, swappable prior has made FM a default choice in large-scale generative pipelines for images~\cite{esser2024sd3}, audio~\cite{vyas2023audiobox}, and proteins~\cite{huguet2024proteinfm}, and it is increasingly applied to physical simulation, where sampling multiple plausible continuations is a better fit than a single deterministic output. This connects directly to the gap identified above: deterministic PDE foundation models such as Poseidon, MPP, and DPOT are trained as single-trajectory regressors, whereas a flow-matching-based model can reuse the same data and encoder but is instead trained to represent a distribution over future states, offering a natural way to test whether explicitly modelling branching behaviour improves long-horizon rollout of Navier-Stokes systems.
\begin{tcolorbox}[
  colback=grey!10,
  colframe=grey!10,
  boxrule=0pt,
  sharp corners,
  enhanced
]
In this paper we introduce NavierNet/FluidGPT, one deterministic and one probabilistic foundation model trained on nine different Navier-Stokes-specific datasets that differ in equation form, parameter settings, and boundary and initial conditions. Both models are trained on a combined pretraining corpus of $\sim$20\,TB, an order of magnitude larger than comparable single-PDE-family foundation models. The deterministic model fuses a Swin Transformer with a classical U-Net backbone and predicts next-timebundle states autoregressively. We compare this model against a probabilistic counterpart based on flow matching, which is designed to capture the branching behaviour of chaotic dynamical systems rather than committing to a single deterministic trajectory. Concretely, we address the following questions:
\begin{enumerate}
    \item Can we develop a rollout-stable, physics-adhering autoregressive deterministic model that generalises across multiple different properties of a Navier-Stokes fluid system using only velocity-field inputs, trained on large-scale pretraining data from different origins?
    \item Can a probabilistic, flow-matching-based approach outperform the deterministic model by better capturing the bifurcative behaviour of high-dimensional Navier-Stokes systems?
\end{enumerate}
\end{tcolorbox}